\hypertarget{ux5e1ux5d9ux5dbux5d5ux5dd-ux5d5ux5e2ux5d1ux5d5ux5d3ux5d4-ux5e2ux5eaux5d9ux5d3ux5d9ux5ea}{%
\section{סיכום ועבודה
עתידית}\label{ux5e1ux5d9ux5dbux5d5ux5dd-ux5d5ux5e2ux5d1ux5d5ux5d3ux5d4-ux5e2ux5eaux5d9ux5d3ux5d9ux5ea}}

\hypertarget{ux5d0ux5eaux5d2ux5e8ux5d9ux5dd-ux5deux5d5ux5e8ux5e4ux5d5ux5dcux5d5ux5d2ux5d9ux5d9ux5dd-ux5e9ux5dc-ux5e2ux5d1ux5e8ux5d9ux5ea}{%
\subsection{אתגרים מורפולוגיים של
עברית}\label{ux5d0ux5eaux5d2ux5e8ux5d9ux5dd-ux5deux5d5ux5e8ux5e4ux5d5ux5dcux5d5ux5d2ux5d9ux5d9ux5dd-ux5e9ux5dc-ux5e2ux5d1ux5e8ux5d9ux5ea}}

עיבוד שפה עברית מציב אתגרים ייחודיים בפני מודלים מבוססי טרנספורמר,
הנובעים מן המורפולוגיה העשירה של השפה.

עברית היא שפה שמית בעלת מורפולוגיה לא-רציפה: שורשים של שלוש או ארבע
אותיות משולבים בתבניות בניין ומשקל כדי לייצר מגוון רחב של צורות מילה.
תהליך זה שונה מהותית מן הלקסיקות הצירופיות המאפיינות שפות
אינדו-אירופיות. כתוצאה מכך, מספר ה-types הייחודיים בטקסט עברי נטוי הוא
גדול בהרבה מזה של שפות כמו אנגלית, מה שמגדיל את מרחב הטוקנים ומקשה על
כיסוי הלקסיקון.

אתגר נוסף הוא הצמדת מילות יחס וצירופי שם: מילות תפקוד כגון ב, כ, ל, מ,
ו, ה מוצמדות ישירות לשם העצם ומהוות עמו יחידת כתיב אחת. הפרדתן דורשת
מנגנוני פיצול שמבצעים ניתוח מורפולוגי בטרם הטוקניזציה. יתרה מכך, בהיעדר
ניקוד בטקסט בלתי מנוקד, קיימת אי-שקיפות פונולוגית המחייבת ביצוע ניתוח
מסמנים לפני ייחוס הפרשנות הנכונה.

\hypertarget{alephbert-ux5d5ux5deux5d5ux5d3ux5dcux5d9ux5dd-ux5e2ux5d1ux5e8ux5d9ux5d9ux5dd}{%
\subsection{AlephBERT ומודלים
עבריים}\label{alephbert-ux5d5ux5deux5d5ux5d3ux5dcux5d9ux5dd-ux5e2ux5d1ux5e8ux5d9ux5d9ux5dd}}

בשנים האחרונות פותחו מודלי שפה מאומנים מראש המותאמים במיוחד לעברית
המודרנית.

מודל AlephBERT אומן על קורפוס עברי נרחב הכולל טקסטים מן האינטרנט, מאמרים
עיתונאיים וספרות, ומשתמש בארכיטקטורת BERT עם טוקניזציה מותאמת המבוססת על
BPE. מחקר זה הדגים כי אימון מראש על קורפוס חד-לשוני עברי עולה בביצועיו
על mBERT בכל המשימות שנבחנו בבנצ'מרק HebrewNLP, לרבות זיהוי ישויות,
ניתוח תחבירי ועוד.

יישום ארכיטקטורת BERT על עברית דרש פתרון סוגיות טוקניזציה ייחודיות.
בניגוד לשפות שבהן WordPiece מייצר פיצולים תחומי-מורפמה, בעברית פיצולים
על-בסיס תדירות בלבד עלולים להפריד יחידות משמעותיות באופן שרירותי. עבודות
מחקר עדכניות בוחנות שילוב של ניתוח מורפולוגי כשלב עיבוד מקדים
לטוקניזציה, במטרה לשפר את איכות הייצוגים הנוצרים.

\hypertarget{ux5d4ux5e2ux5d1ux5e8ux5d4-ux5d7ux5d5ux5e6ux5ea-ux5e9ux5e4ux5d5ux5ea-ux5d5ux5dbux5d9ux5d5ux5d5ux5e0ux5d9ux5dd-ux5e2ux5eaux5d9ux5d3ux5d9ux5d9ux5dd}{%
\subsection{העברה חוצת-שפות וכיוונים
עתידיים}\label{ux5d4ux5e2ux5d1ux5e8ux5d4-ux5d7ux5d5ux5e6ux5ea-ux5e9ux5e4ux5d5ux5ea-ux5d5ux5dbux5d9ux5d5ux5d5ux5e0ux5d9ux5dd-ux5e2ux5eaux5d9ux5d3ux5d9ux5d9ux5dd}}

מודל mBERT, שאומן על ויקיפדיה ב-104 שפות, הפגין יכולות העברה אפסית
מרשימות בין שפות, כולל עברית. עם זאת, הפערים ביחס למודלים חד-לשוניים
ייעודיים מצביעים על כך שהעברה חוצת-שפות אינה תחליף מלא לאימון על קורפוס
ייעודי.

כיוונים עתידיים בתחום כוללים: פיתוח ארכיטקטורות Transformer יעילות יותר
המצמצמות את מורכבות O(n\^{}2) של מנגנון הקשב; בניית קורפוסים עבריים
מאוגדים וגדולים יותר; שילוב מידע מורפולוגי מפורש בשכבות הייצוג; ומחקר
בתחום ה-LLM הרב-לשוני.

\hypertarget{ux5e1ux5d9ux5dbux5d5ux5dd-ux5dbux5dcux5dcux5d9}{%
\subsection{סיכום
כללי}\label{ux5e1ux5d9ux5dbux5d5ux5dd-ux5dbux5dcux5dcux5d9}}

פרק זה סקר את האתגרים הייחודיים של עיבוד שפה עברית במסגרת ארכיטקטורת
הטרנספורמר, ואת הפתרונות שפותחו לאחרונה. עבודה עתידית בתחום תידרש להמשיך
ולגשר בין ההתקדמות הכללית בתחום מודלי השפה הגדולים לבין הצרכים הייחודיים
של שפות מורפולוגית מורכבות כגון עברית.
